% संपूर्ण मराठी ग्रंथातील प्रकरण 12: स्वयंसिद्धकीय निष्पत्ती
% हा संपादनयोग्य स्रोतखंड आहे. संकलनासाठी संपूर्ण स्रोतसंग्रहातील
% अक्षररूपे, पूर्वभाग, चिन्हव्याख्या आणि आकृती-स्रोत आवश्यक आहेत.
% मूळ: Open Logic Project; मजकूर आणि रूपांतर CC BY 4.0.
% यंत्रानुवाद, दुरुस्ती आणि तपासणी: OpenAI Codex —
% GPT-5.6 Sol आणि GPT-6 Sol, दोन्ही Ultra effort.
% दुरुस्ती-पुनरावलोकन: GPT-6.1 Sol, Ultra effort.
\chapter{स्वयंसिद्धकीय निष्पत्ती}\label{fol:axd::chap}
\begin{editorial}
  या प्रकरणातील मजकूर कोणते संयोजक आणि संख्यापक आदिम किंवा परिभाषित
  आहेत हे दर्शवणाऱ्या विविध टॅग्जचे पालन करतो याची खात्री करण्याचा अद्याप
  कोणताही प्रयत्न केलेला नाही: यांपैकी $\liff$ परिभाषित आहे, तर बाकी
  सर्व आदिम आहेत असे गृहीत धरले आहे. FOL टॅग सत्य
  असल्यास आपण संख्यापकांसह आवृत्ती तयार करतो; अन्यथा संख्यापकांविना
  आवृत्ती तयार करतो.
\end{editorial}









\section{नियम आणि निष्पत्ती}\label{fol:axd:rul:sec}

\begin{explain}
  स्वयंसिद्धकीय निष्पत्त्या ही तर्कशास्त्रासाठी कदाचित सर्वांत सोपी
  निष्पत्ती-पद्धत आहे. निष्पत्ती म्हणजे केवळ सूत्रांचा
  अनुक्रम. निष्पत्ती म्हणून गणले जाण्यासाठी अनुक्रमातील प्रत्येक
  सूत्र एकतर स्वयंसिद्धकाचे उदाहरण असले पाहिजे, किंवा अनुक्रमात
  त्याच्या आधी येणाऱ्या एक वा अधिक सूत्रांपासून निगमन नियमाने
  निष्पन्न झाले पाहिजे. निष्पत्ती तिचे शेवटचे सूत्र
  निष्पन्न करते.
\end{explain}

\begin{defn}[निष्पन्नता]
जर $\Gamma$ हा $\Lang L$ मधील सूत्रांचा संच असेल, तर
$\Gamma$ पासूनची \emph{निष्पत्ती} म्हणजे सूत्रांचा सांत अनुक्रम
$\varphi_1$, \dots,~$\varphi_n$ असा की प्रत्येक $i \le n$ साठी पुढीलपैकी एक अट
पूर्ण होते:
\begin{enumerate}
\item $\varphi_i \in \Gamma$; किंवा
\item $\varphi_i$ हे स्वयंसिद्धक आहे; किंवा
\item $j < i$ (आणि $k < i$) असलेल्या एखाद्या $\varphi_j$ (आणि $\varphi_k$)
  पासून निगमन नियमाने $\varphi_i$ निष्पन्न होते.
\end{enumerate}
\end{defn}

योग्य निष्पत्ती म्हणून कशाची गणना होते हे आपण कोणते निगमन नियम
मान्य करतो (आणि अर्थातच कोणती सूत्रे स्वयंसिद्धके मानतो) यावर अवलंबून
असते. निगमन नियम म्हणजे असे जर-तर विधान जे विशिष्ट अटींमध्ये
निष्पत्तीमधील पायरी~$A_i$ ही निगमनाची योग्य पायरी आहे असे सांगते.

\begin{defn}[निगमन नियम]
\emph{निगमन नियम} हा $\Gamma$ पासूनच्या निष्पत्तीमध्ये निगमनाची
योग्य पायरी म्हणून कशाची गणना होते यासाठी पुरेशी अट देतो.
\end{defn}

उदाहरणार्थ, $\varphi \in \Gamma$ असलेला कोणताही एक-घटकी अनुक्रम $\varphi$ हा
आपोआपच निष्पत्ती म्हणून गणला जात असल्यामुळे, पुढील नियम हा
अतिशय सोपा निगमन नियम असू शकतो:
\begin{quote}
  जर $\varphi \in \Gamma$, तर $\Gamma$ पासूनच्या कोणत्याही निष्पत्तीमध्ये $\varphi$ ही नेहमीच निगमनाची योग्य पायरी असते.
\end{quote}
त्याचप्रमाणे, $\varphi$ हे स्वयंसिद्धकांपैकी एक असेल, तर एकटे $\varphi$ ही
निष्पत्ती असते; म्हणून पुढील विधानही निगमन नियम आहे:
\begin{quote}
  जर $\varphi$ हे स्वयंसिद्धक असेल, तर $\varphi$ ही निगमनाची योग्य पायरी असते.
\end{quote}
विचाराधीन पायरीच्या आधी दिसणाऱ्या सूत्रांचा निगमन नियम आधार घेतो
तेव्हा गोष्ट अधिक रोचक होते. पुढील नियमाला \emph{विध्यात्मक अनुमान}
म्हणतात:
\begin{quote}
  जर $\psi \lif \varphi$ आणि $\psi$ हे निष्पत्तीमध्ये यापूर्वी आढळत असतील,
  तर~$\varphi$ ही निगमनाची योग्य पायरी असते.
\end{quote}
हा एकमेव निगमन नियम असेल, तर वर दिलेल्या निष्पत्तीच्या व्याख्येचा
अर्थ असा होतो: प्रत्येक $i \le n$ साठी पुढीलपैकी एक अट पूर्ण होत असेल
तेव्हा आणि केवळ तेव्हाच $\varphi_1$, \dots,~$\varphi_n$ ही निष्पत्ती असते:
\begin{enumerate}
\item $\varphi_i \in \Gamma$; किंवा
\item $\varphi_i$ हे स्वयंसिद्धक आहे; किंवा
\item एखाद्या $j < i$ साठी $\varphi_j$ हे $\psi \lif \varphi_i$ आहे, आणि एखाद्या
  $k < i$ साठी $\varphi_k$ हे~$\psi$ आहे.
\end{enumerate}
शेवटचे कलम असे सांगते की~$\varphi_j$ ($\psi \lif \varphi_i$) आणि $\varphi_k$ ($\psi$)
यांपासून विध्यात्मक अनुमानाने $\varphi_i$ निष्पन्न होते. आपण $1$ पासून~$n$
पर्यंत जाऊ शकलो आणि प्रत्येक वेळी~$\Gamma$ मध्ये असलेले, स्वयंसिद्धक
असलेले किंवा निगमन नियमाने योग्य निगमन-पायरी ठरणारे सूत्र~$\varphi_i$
मिळाले, तर संपूर्ण अनुक्रम योग्य निष्पत्ती म्हणून गणला जातो.

\begin{defn}[निष्पन्नता]
जर $\Gamma$ पासून सुरू होऊन $\varphi$ वर संपणारी निष्पत्ती असेल, तर
सूत्र~$\varphi$ हे $\Gamma$ पासून \emph{निष्पन्न करता येण्याजोगे} असते; हे आपण
$\Gamma \Proves \varphi$ असे लिहितो.
\end{defn}

\begin{defn}[प्रमेये]
रिक्त संचापासून $\varphi$ ची निष्पत्ती असेल, तर सूत्र~$\varphi$ हे
\emph{प्रमेय} असते. $\varphi$ हे प्रमेय असेल तर आपण $\Proves \varphi$ लिहितो आणि
ते प्रमेय नसेल तर $\Proves/ \varphi$ लिहितो.
\end{defn}













\section{विधानीय संयोजकांसाठी स्वयंसिद्धके आणि नियम}\label{fol:axd:prp:sec}

\begin{defn}[स्वयंसिद्धके]
विधानीय संयोजकांसाठीच्या \emph{स्वयंसिद्धकांचा} संच $\PAx$ हा पुढील रूपांच्या
सर्व सूत्रांचा बनलेला आहे:
\begin{align}
  & (\varphi \land \psi) \lif \varphi \label{fol:axd:prp:ax:land1}\\
  & (\varphi \land \psi) \lif \psi \label{fol:axd:prp:ax:land2}\\
  & \varphi \lif (\psi \lif (\varphi \land \psi)) \label{fol:axd:prp:ax:land3}\\
  & \varphi \lif (\varphi \lor \psi) \label{fol:axd:prp:ax:lor1}\\
  & \varphi \lif (\psi \lor \varphi) \label{fol:axd:prp:ax:lor2}\\
  & (\varphi \lif \chi) \lif ((\psi \lif \chi) \lif ((\varphi \lor \psi) \lif \chi)) \label{fol:axd:prp:ax:lor3}\\
  & \varphi \lif (\psi \lif \varphi) \label{fol:axd:prp:ax:lif1}\\
  & (\varphi \lif (\psi \lif \chi)) \lif ((\varphi \lif \psi) \lif (\varphi \lif \chi)) \label{fol:axd:prp:ax:lif2}\\
  & (\varphi \lif \psi) \lif ((\varphi \lif \lnot \psi) \lif \lnot \varphi) \label{fol:axd:prp:ax:lnot1}\\
  & \lnot \varphi \lif (\varphi \lif \psi) \label{fol:axd:prp:ax:lnot2}\\
  & \ltrue \label{fol:axd:prp:ax:ltrue}\\
  & \lfalse \lif \varphi \label{fol:axd:prp:ax:lfalse1}\\
  & (\varphi \lif \lfalse) \lif \lnot \varphi \label{fol:axd:prp:ax:lfalse2}\\
  & \lnot\lnot \varphi \lif \varphi \label{fol:axd:prp:ax:dne}
\end{align}
\end{defn}

\begin{defn}[विध्यात्मक अनुमान]
  जर $\psi$ आणि $\psi \lif \varphi$ हे निष्पत्तीमध्ये आधीच आलेले असतील,
  तर $\varphi$ ही निगमनाची योग्य पायरी असते.
\end{defn}

विध्यात्मक अनुमान या नियमाचा संक्षेप आपण ``\MP'' असा करू.












\section{संख्यापकांसाठी स्वयंसिद्धके आणि नियम}\label{fol:axd:qua:sec}

\begin{defn}[संख्यापकांसाठी स्वयंसिद्धके]
संख्यापकांचे नियमन करणारी \emph{स्वयंसिद्धके} म्हणजे पुढील रूपांचे सर्व
दृष्टांत होत:
\begin{align}
\label{fol:axd:qua:ax:q1} & \lforall[x][\psi] \lif \psi(t), \\
\label{fol:axd:qua:ax:q2} & \psi(t) \lif \lexists[x][\psi].
\end{align}
कोणत्याही बंद पद~$t$ साठी.
\end{defn}

\begin{defn}[संख्यापकांसाठी नियम]
\item जर $\psi \lif \varphi(a)$ हे निष्पत्तीमध्ये आधीच आलेले असेल आणि
  $a$ ची~$\Gamma$ किंवा~$\psi$ मध्ये आस्थिति नसेल, तर $\psi \lif
  \lforall[x][\varphi(x)]$ ही निगमनाची योग्य पायरी असते.
\item जर $\varphi(a) \lif \psi$ हे निष्पत्तीमध्ये आधीच आलेले असेल आणि
  $a$ ची~$\Gamma$ किंवा~$\psi$ मध्ये आस्थिति नसेल, तर
  $\lexists[x][\varphi(x)] \lif \psi$ ही निगमनाची योग्य पायरी असते.
\end{defn}

यांपैकी कोणत्याही नियमाचा संक्षेप आपण ``\QR'' असा करू.












\section{निष्पत्ती: उदाहरणे}\label{fol:axd:pro:sec}


\begin{ex}
आपल्याला $(\lnot \theta \lor \varphi_{E}) \lif (\theta \lif
\varphi_{E})$ हे सिद्ध करायचे आहे असे समजा. स्पष्टपणे, हे आपल्या कोणत्याही
स्वयंसिद्धकाचे उदाहरण नाही; म्हणून ते निष्पन्न करण्यासाठी \MP{} नियम
वापरावा लागेल. आपला एकमेव नियम~MP आहे; $\varphi$ आणि $\varphi \lif \psi$ दिले असता
या नियमाने~$\psi$ ला समर्थित करता येते. एक युक्ती अशी असेल: $\varphi$ म्हणून
$\lnot \theta$, $\psi$ म्हणून $\varphi_{E}$ आणि $\chi$ म्हणून $\theta \lif \varphi_{E}$ घेऊन
\ref{fol:axd:prp:ax:lor3} वापरणे, म्हणजे पुढील उदाहरण वापरणे:
\[(\lnot \theta \lif (\theta \lif \varphi_{E})) \lif ((\varphi_{E} \lif (\theta \lif \varphi_{E})) \lif ((\lnot
\theta \lor \varphi_{E}) \lif (\theta \lif \varphi_{E}))).
\]
का? MP चे दोन उपयोग केल्यावर शेवटचा भाग मिळतो, आणि आपल्याला नेमके तेच
हवे आहे. तसेच $\lnot \theta \lif (\theta \lif \varphi_{E})$ हे
\ref{fol:axd:prp:ax:lnot2} चे उदाहरण आणि $\varphi_{E} \lif (\theta \lif \varphi_{E})$ हे
\ref{fol:axd:prp:ax:lif1} चे उदाहरण आहे, हे सहज दिसते. म्हणून आपली
निष्पत्ती अशी आहे:
\begin{derivation}
  1. &  $\lnot \theta \lif (\theta \lif \varphi_{E})$ & \ref{fol:axd:prp:ax:lnot2} \\
  2. & $(\lnot \theta \lif (\theta \lif \varphi_{E})) \lif {}$\\
  &\qquad $((\varphi_{E} \lif (\theta \lif \varphi_{E})) \lif ((\lnot \theta \lor \varphi_{E}) \lif (\theta \lif \varphi_{E})))$ & \ref{fol:axd:prp:ax:lor3}\\
  3. & $(\varphi_{E} \lif (\theta \lif \varphi_{E})) \lif ((\lnot \theta \lor \varphi_{E}) \lif (\theta \lif \varphi_{E}))$ & 1, 2, \MP\\
  4. & $\varphi_{E} \lif (\theta \lif \varphi_{E})$ & \ref{fol:axd:prp:ax:lif1}\\
  5. & $(\lnot \theta \lor \varphi_{E}) \lif (\theta \lif \varphi_{E})$ & 3, 4, \MP
\end{derivation}
\end{ex}

\begin{ex}\label{fol:axd:pro:ex:identity}
$\theta \lif \theta$ ची निष्पत्ती शोधून पाहू. हे स्वयंसिद्धकाचे उदाहरण
नाही; म्हणून ते निष्पन्न करण्यासाठी आपल्याला \MP{} वापरावा लागेल.
\ref{fol:axd:prp:ax:lif1} हे~$\varphi \lif \psi$ या रूपाचे स्वयंसिद्धक आहे आणि त्याला
आपण~\MP लावू शकतो. त्याचा उपयोग व्हायचा असेल, तर अर्थातच या वेळी
\MP{} ने योग्य पायरी म्हणून समर्थित होणारे $\psi$ हे~$\theta \lif \theta$ असले
पाहिजे, कारण आपल्याला तेच निष्पन्न करायचे आहे. त्यामुळे $\varphi$ सुद्धा
$\theta$ असले पाहिजे; म्हणजे आपण \ref{fol:axd:prp:ax:lif1} चे पुढील उदाहरण पाहू शकतो:
\[\theta \lif (\theta \lif \theta)
\]
\MP लावण्यासाठी आपल्याला संबंधित दुसरे आधारविधान, म्हणजे~$\varphi$, सुद्धा
समर्थित करावे लागेल. पण आपल्या प्रकरणात ते~$\theta$ असेल, आणि एकटे~$\theta$
निष्पन्न करता येणार नाही. म्हणून आपल्याला वेगळी युक्ती हवी.

केवळ~$\lif$ असलेले दुसरे स्वयंसिद्धक म्हणजे \ref{fol:axd:prp:ax:lif2}, म्हणजे
\[(\varphi \lif (\psi \lif \chi)) \lif ((\varphi \lif \psi) \lif (\varphi \lif \chi))
\]
\MP{} दोनदा लावून आपण शेवटच्या अंतःस्थ सशर्त विधानापर्यंत पोहोचू शकतो.
पुन्हा, याचा अर्थ असा की \ref{fol:axd:prp:ax:lif2} चे असे उदाहरण आपल्याला हवे
ज्यात $\varphi \lif \chi$ हे $\theta \lif \theta$, म्हणजे आपण साध्य करू पाहत असलेले
सूत्र, असेल. मग अर्थातच $\varphi$ आणि $\chi$ ही दोन्ही~$\theta$ असतात.
$\varphi \lif (\psi \lif \chi)$ आणि $\varphi \lif \psi$, म्हणजे आपल्या प्रकरणात
$\theta \lif (\psi \lif \theta)$ आणि $\theta \lif \psi$, ही दोन्हीही निष्पन्न करता येण्याजोगे
होतील अशा रीतीने~$\psi$ कसे निवडावे? आपण $\psi$ काहीही निवडले, तरी यांपैकी
पहिले हे आधीच \ref{fol:axd:prp:ax:lif1} चे उदाहरण आहे. आणि $\psi$ हे
$(\theta \lif \theta)$ असेल, तर $\theta \lif \psi$ हे \ref{fol:axd:prp:ax:lif1} चे आणखी
एक उदाहरण असेल. म्हणून आपली निष्पत्ती अशी आहे:
\begin{derivation}
  1. & $\theta \lif ((\theta \lif \theta) \lif \theta)$ & \ref{fol:axd:prp:ax:lif1}\\
  2. & $(\theta \lif ((\theta \lif \theta) \lif \theta)) \lif {}$\\
  & \qquad $((\theta \lif (\theta \lif \theta)) \lif (\theta \lif \theta))$ & \ref{fol:axd:prp:ax:lif2}\\
  3. & $(\theta \lif (\theta \lif \theta)) \lif (\theta \lif \theta)$ & 1, 2, \MP\\
  4. & $\theta \lif (\theta \lif \theta)$ & \ref{fol:axd:prp:ax:lif1}\\
  5. & $\theta \lif \theta$ & 3, 4, \MP
\end{derivation}
\end{ex}

\begin{ex}\label{fol:axd:pro:ex:chain}
कधी कधी काही इतर सूत्रे~$\Gamma$ पासून एखाद्या सूत्राची
निष्पत्ती आहे हे आपल्याला दाखवायचे असते. उदाहरणार्थ,
$\Gamma = \{\varphi \lif \psi, \psi \lif \chi\}$ पासून $\varphi \lif \chi$ हे
निष्पन्न करता येते, हे दाखवू.
\begin{derivation}
  1. & $\varphi \lif \psi$ & \Hyp\\
  2. & $\psi \lif \chi$ & \Hyp\\
  3. & $(\psi \lif \chi) \lif (\varphi \lif (\psi \lif \chi))$ & \ref{fol:axd:prp:ax:lif1} \\
  4. & $\varphi \lif (\psi \lif \chi)$ & 2, 3, \MP\\
  5. & $(\varphi \lif (\psi \lif \chi)) \lif {}$\\
  & \qquad $((\varphi \lif \psi) \lif (\varphi \lif \chi))$ & \ref{fol:axd:prp:ax:lif2}\\
  6. & $((\varphi \lif \psi) \lif (\varphi \lif \chi))$ & 4, 5, \MP\\
  7. & $\varphi \lif \chi$ & 1, 6, \MP
\end{derivation}
``\Hyp'' (``hypothesis'', म्हणजे ``गृहीतक'') असे लेबल असलेल्या ओळी
त्या ओळीवरील सूत्र हा~$\Gamma$ चा घटक असल्याचे दर्शवतात.
\end{ex}

\begin{prop}
  \label{fol:axd:pro:prop:chain} जर $\Gamma \Proves \varphi \lif \psi$ आणि $\Gamma
  \Proves \psi \lif \chi$, तर $\Gamma \Proves \varphi \lif \chi$
\end{prop}

\begin{proof}
  $\Gamma \Proves \varphi \lif \psi$ आणि $\Gamma \Proves \psi \lif
  \chi$ असे समजा. मग~$\Gamma$ पासून $\varphi \lif \psi$ ची निष्पत्ती
  असते; तसेच~$\Gamma$ पासून $\psi \lif \chi$ चीही निष्पत्ती असते.
  त्यांची जोडणी करून त्या एकाच निष्पत्तीमध्ये एकत्र करा. आता मागील
  उदाहरणातील निष्पत्तीच्या ओळी 3--7 जोडा. ही~$\Gamma$ पासूनची
  निष्पत्ती आहे, आणि नव्या निष्पत्तीची शेवटची ओळ
  $\varphi \lif \chi$ आहे. लक्षात घ्या की ओळी 4 आणि~7 साठी दिलेली
  समर्थने पुढील बदलांनंतरही वैध राहतात: ओळ क्रमांक~1 चा निर्देश $\varphi
  \lif \psi$ च्या निष्पत्तीच्या शेवटच्या ओळीच्या निर्देशाने आणि ओळ
  क्रमांक~2 चा निर्देश $\psi \lif \chi$ च्या निष्पत्तीच्या शेवटच्या
  ओळीच्या निर्देशाने बदला.
\end{proof}

\begin{prob}
स्वयंसिद्धकांपासूनच्या निष्पत्त्या मांडून पुढील विधाने सिद्ध करा:
\begin{enumerate}
\item $(\varphi \land \psi) \lif (\psi \land \varphi)$
\item $((\varphi \land \psi) \lif \chi) \lif (\varphi \lif (\psi \lif \chi))$
\item $\lnot(\varphi \lor \psi) \lif \lnot \varphi$
\end{enumerate}
\end{prob}














\section{संख्यापकांसह निष्पत्ती}\label{fol:axd:prq:sec}

\begin{ex}
$(\lforall[x][\varphi(x)] \land
\lforall[y][\psi(y)]) \lif \lforall[x][(\varphi(x) \land \psi(x))]$ ची
निष्पत्ती देऊ.

प्रथम, लक्षात घ्या की
\begin{align*}
  (\lforall[x][\varphi(x)] \land \lforall[y][\psi(y)]) & \lif \lforall[x][\varphi(x)]\\
  \intertext{हे \ref{fol:axd:prp:ax:land1} चे उदाहरण आहे, आणि}
  \lforall[x][\varphi(x)] & \lif \varphi(a) \\
  \intertext{हे \ref{fol:axd:qua:ax:q1} चे उदाहरण आहे. म्हणून \ref{fol:axd:pro:prop:chain} वरून आपल्याला माहीत आहे की}
  (\lforall[x][\varphi(x)] \land \lforall[y][\psi(y)]) & \lif \varphi(a) \\
  \intertext{हे निष्पन्न करता येण्याजोगे आहे. त्याचप्रमाणे,}
  (\lforall[x][\varphi(x)] \land \lforall[y][\psi(y)]) & \lif \lforall[y][\psi(y)] \qquad\text{आणि}\\
    \lforall[y][\psi(y)] & \lif \psi(a)\\
    \intertext{ही अनुक्रमे \ref{fol:axd:prp:ax:land2} आणि \ref{fol:axd:qua:ax:q1} यांची उदाहरणे असल्यामुळे,}
    (\lforall[x][\varphi(x)] \land \lforall[y][\psi(y)]) & \lif \psi(a)\\
    \intertext{हे \ref{fol:axd:pro:prop:chain} मुळे सिद्धतायोग्य आहे. \ref{fol:axd:prp:ax:land3} चे योग्य उदाहरण आणि~\MP चे दोन उपयोग करून आपल्याला दिसते की}
    (\lforall[x][\varphi(x)] \land \lforall[y][\psi(y)]) & \lif (\varphi(a) \land \psi(a))\\
    \intertext{हे सिद्धतायोग्य आहे. आता \QR{} लावून पुढील सूत्र मिळते:}
(\lforall[x][\varphi(x)] \land \lforall[y][\psi(y)]) & \lif \lforall[x][(\varphi(x) \land \psi(x))].
\end{align*}
\end{ex}












\section{सिद्धता-उपपत्तीय संकल्पना}\label{fol:axd:ptn:sec}

\begin{explain}
जशा आपण अनेक महत्त्वाच्या चिन्हार्थविषयक संकल्पना
(वैधता, तार्किक निष्पन्नता, पूर्ततायोग्यता)
परिभाषित केल्या आहेत, तशाच आता त्यांना अनुरूप
\emph{सिद्धता-उपपत्तीय संकल्पना} परिभाषित करू. या संकल्पना संरचना
मध्ये वाक्यांची पूर्ति होते की नाही याचा आधार घेऊन परिभाषित केल्या
जात नाहीत; त्याऐवजी विशिष्ट सूत्रांची निष्पन्नता किंवा
अनिष्पन्नता यांचा आधार घेतला जातो. या दोन प्रकारच्या संकल्पना
एकमेकांशी जुळतात, हा एक महत्त्वाचा शोध होता. त्या जुळतात, हाच
\emph{निर्दोषता प्रमेय} आणि \emph{संपूर्णता प्रमेय} यांचा आशय आहे.
\end{explain}

\begin{defn}[निष्पन्नता]
सूत्र~$\varphi$ हे $\Gamma$ \emph{पासून निष्पन्न करता येण्याजोगे} असते, म्हणजे
$\Gamma \Proves \varphi$, जर~$\Gamma$ पासून सुरू होऊन~$\varphi$ वर संपणारी
निष्पत्ती असेल.
\end{defn}

\begin{defn}[प्रमेये]
रिक्त संचापासून $\varphi$ ची निष्पत्ती असेल, तर सूत्र~$\varphi$ हे
\emph{प्रमेय} असते. $\varphi$ हे प्रमेय असेल तर आपण $\Proves \varphi$ लिहितो आणि
ते प्रमेय नसेल तर $\Proves/ \varphi$ लिहितो.
\end{defn}

\begin{defn}[सुसंगतता]
सूत्रांचा संच~$\Gamma$ हा $\Gamma\Proves/ \lfalse$ असेल तेव्हा आणि
केवळ तेव्हाच \emph{सुसंगत} असतो; अन्यथा तो \emph{विसंगत} असतो.
\end{defn}

\begin{prop}[परावर्तिता]
\label{fol:axd:ptn:prop:reflexivity}
$\varphi \in \Gamma$ असेल, तर $\Gamma \Proves \varphi$.
\end{prop}

\begin{proof}
  एकटे सूत्र~$\varphi$ ही~$\Gamma$ पासून $\varphi$ ची निष्पत्ती आहे.
\end{proof}

\begin{prop}[एकस्वनिकता]
\label{fol:axd:ptn:prop:monotonicity}
$\Gamma \subseteq \Delta$ आणि $\Gamma \Proves \varphi$ असेल, तर $\Delta
\Proves \varphi$.
\end{prop}

\begin{proof}
$\Gamma$ पासून $\varphi$ ची कोणतीही निष्पत्ती ही~$\Delta$ पासून $\varphi$ चीही
निष्पत्ती असते.
\end{proof}

\begin{prop}[संक्रमकता]
\label{fol:axd:ptn:prop:transitivity}
$\Gamma \Proves \varphi$ आणि $\{\varphi\} \cup \Delta \Proves
\psi$ असेल, तर $\Gamma \cup \Delta \Proves \psi$.
\end{prop}

\begin{proof}
  $\{\varphi\} \cup \Delta \Proves \psi$ असे समजा. मग~$\{\varphi\} \cup
  \Delta$ पासून $\psi_1$, \dots, $\psi_l = \psi$ ही निष्पत्ती असते.
  त्या निष्पत्तीमधील काही पायऱ्या योग्य असतात, कारण एखादा नियम
  त्यांपूर्वीच्या~$\psi_i = \varphi$ या ओळीचा निर्देश करतो. गृहीतकानुसार,
  $\Gamma$ पासून~$\varphi$ ची निष्पत्ती असते, म्हणजे निष्पत्ती
  $\varphi_1$, \dots, $\varphi_k = \varphi$ अशी असते की प्रत्येक $\varphi_i$ हे स्वयंसिद्धक,
  $\Gamma$ चा घटक, किंवा निगमन नियमाने योग्य ठरलेले असते. आता
  पुढील अनुक्रम विचारात घ्या:
  \[  \varphi_1, \dots, \varphi_k = \varphi, \psi_1, \dots, \psi_l = \psi.
  \]
  ही $\Gamma \cup \Delta$ पासून~$\psi$ ची योग्य निष्पत्ती आहे, कारण
  प्रत्येक $B_i = \varphi$ ला आता~$\varphi_k = \varphi$ ला समर्थित करणाऱ्या त्याच नियमाने
  समर्थन मिळते.
\end{proof}

लक्षात घ्या की विशेषतः याचा अर्थ असा होतो: $\Gamma \Proves \varphi$ आणि $\varphi
\Proves \psi$ असेल, तर $\Gamma \Proves \psi$. तसेच $\varphi_1,
\dots, \varphi_n \Proves \psi$ असेल आणि प्रत्येक~$i$ साठी $\Gamma \Proves \varphi_i$
असेल, तर $\Gamma \Proves \psi$, हेसुद्धा यावरून मिळते.

\begin{prop}
\label{fol:axd:ptn:prop:incons}
$\Gamma$ विसंगत असतो तेव्हा आणि केवळ तेव्हाच, जेव्हा प्रत्येक~$\varphi$ साठी
$\Gamma \Proves \varphi$.
\end{prop}

\begin{proof}
सराव.
\end{proof}


\begin{prob}
\ref{fol:axd:ptn:prop:incons} सिद्ध करा.
\end{prob}




\begin{prop}[संहतता]
\label{fol:axd:ptn:prop:proves-compact}
  \begin{enumerate}
  \item $\Gamma \Proves \varphi$ असेल, तर एखादा सांत उपसंच $\Gamma_0
    \subseteq \Gamma$ असा असतो की $\Gamma_0 \Proves \varphi$.
  \item $\Gamma$ चा प्रत्येक सांत उपसंच सुसंगत असेल, तर $\Gamma$ सुसंगत
    असतो.
  \end{enumerate}
\end{prop}

\begin{proof}
  \begin{enumerate}
    \item $\Gamma \Proves \varphi$ असेल, तर सूत्रांचा सांत अनुक्रम
      $\varphi_1$, \dots,~$\varphi_n$ असा असतो की $\varphi \ident \varphi_n$ आणि प्रत्येक
      $\varphi_i$ हे एकतर तर्कशास्त्रीय स्वयंसिद्धक, $\Gamma$ चा घटक,
      किंवा आधीच्या सूत्रांपासून विध्यात्मक अनुमानाने निष्पन्न झालेले
      असते. $\Gamma$ मध्ये असलेल्या त्या $\varphi_i$ चा संच~$\Gamma_0$ म्हणून
      घ्या. मग ही निष्पत्ती त्याचप्रमाणे~$\Gamma_0$ पासूनची
      निष्पत्ती असते, आणि म्हणून $\Gamma_0 \Proves \varphi$.
    \item (1) मध्ये $\varphi
      \ident \lfalse$ हे विशेष प्रकरण घेतल्यावर मिळणारे हे निषेधात्मक उलट
      विधान आहे.
\end{enumerate}
\end{proof}















\section{निगमन प्रमेय}\label{fol:axd:ded:sec}

आपण पाहिल्याप्रमाणे, स्वयंसिद्धकीय पद्धतीत निष्पत्त्या देणे किचकट
असते आणि निष्पत्त्या शोधणे कठीण असू शकते. सूत्रांच्या लांब
याद्या प्रत्यक्ष लिहिण्याऐवजी काही सोपे निष्कर्ष वापरून अशा
निष्पत्त्या अस्तित्वात आहेत असा युक्तिवाद करणे सहसा सोपे असते. असे तीन
निष्कर्ष आपण आधीच प्रस्थापित केले आहेत: $\varphi \in \Gamma$ हे माहीत असेल,
तर $\Gamma \Proves \varphi$ असे नेहमी म्हणता येते, हे
\ref{fol:axd:ptn:prop:reflexivity} सांगतो. $\Gamma \Proves \varphi$ असेल, तर
$\Gamma \cup \{\psi\} \Proves \varphi$ सुद्धा असते, हे
\ref{fol:axd:ptn:prop:monotonicity} सांगतो. आणि $\Gamma \Proves \varphi$ व
$\varphi \Proves \psi$ असेल, तर $\Gamma \Proves \psi$, असे
\ref{fol:axd:ptn:prop:transitivity} वरून मिळते. पुढे आणखी एक सोपा निष्कर्ष,
विध्यात्मक अनुमानाची ``अधि''-आवृत्ती, दिली आहे:

\begin{prop}
  \label{fol:axd:ded:prop:mp} जर $\Gamma \Proves \varphi$ आणि $\Gamma \Proves \varphi \lif
  \psi$, तर $\Gamma \Proves \psi$.
\end{prop}

\begin{proof}
  $\{\varphi, \varphi \lif \psi\} \Proves \psi$ हे आपल्याकडे आहे:
  \begin{derivation}
    1. & $\varphi$ & Hyp.\\
    2. & $\varphi \lif \psi$ & Hyp.\\
    3. & $\psi$ & 1, 2, MP
  \end{derivation}
  \ref{fol:axd:ptn:prop:transitivity} वरून $\Gamma \Proves \psi$.
\end{proof}

या संदर्भात आपण वापरणार असलेला सर्वांत महत्त्वाचा निष्कर्ष म्हणजे निगमन
प्रमेय:

\begin{thm}[निगमन प्रमेय]
\label{fol:axd:ded:thm:deduction-thm} $\Gamma \cup \{\varphi\} \Proves \psi$ तेव्हा आणि
केवळ तेव्हाच, जेव्हा $\Gamma \Proves \varphi \lif \psi$.
\end{thm}

\begin{proof}
``जर'' ही दिशा तात्काळ मिळते. $\Gamma \Proves \varphi \lif \psi$ असेल, तर
\ref{fol:axd:ptn:prop:monotonicity} वरून $\Gamma \cup \{\varphi\} \Proves \varphi \lif \psi$
सुद्धा असते. तसेच \ref{fol:axd:ptn:prop:reflexivity} वरून
$\Gamma \cup \{\varphi\} \Proves \varphi$. म्हणून \ref{fol:axd:ded:prop:mp} वरून
$\Gamma \cup \{\varphi\} \Proves \psi$.

``केवळ जर'' या दिशेसाठी, $\Gamma \cup \{\varphi\}$ पासून $\psi$ च्या
निष्पत्तीच्या लांबीवर विगमन करू.

विगमनाच्या आधारपायरीसाठी, लांबी~$1$ असलेल्या प्रत्येक निष्पत्ती साठी
दावा सिद्ध करू. $\Gamma \cup \{\varphi\}$ पासून~$\psi$ ची लांबी~$1$ असलेली
निष्पत्ती म्हणजे एकटे $\psi$; ती योग्य असेल, तर $\psi$ हे एकतर
$\in \Gamma \cup \{\varphi\}$ असते किंवा स्वयंसिद्धक असते. $\psi \in \Gamma$
असेल किंवा ते स्वयंसिद्धक असेल, तर $\Gamma \Proves \psi$. तसेच
\ref{fol:axd:prp:ax:lif1} वरून $\Gamma \Proves \psi \lif (\varphi \lif \psi)$, आणि
\ref{fol:axd:ded:prop:mp} वरून $\Gamma \Proves \varphi \lif \psi$. $\psi \in \{ \varphi\}$
असेल, तर $\Gamma \Proves \varphi \lif \psi$, कारण शेवटचे वाक्य~$\varphi
\lif \psi$ हे $\varphi \lif \varphi$ सारखेच आहे आणि ते आपण
\ref{fol:axd:pro:ex:identity} मध्ये निष्पन्न केले आहे.

विगमन पायरीसाठी, $\Gamma \cup \{\varphi\}$ पासून~$\psi$ ची एखादी
निष्पत्ती विध्यात्मक अनुमानाने समर्थित असलेल्या~$\psi$ या पायरीने
संपते असे समजा. (तिला विध्यात्मक अनुमानाने समर्थन मिळत नसेल, तर
$\psi \in \Gamma$, $\psi \ident \varphi$, किंवा $\psi$ हे स्वयंसिद्धक असते, आणि
विगमनाच्या आधारपायरीतील युक्तिवादच लागू होतो.) मग एखाद्या
सूत्र~$\chi$ साठी त्या निष्पत्तीमधील आधीच्या काही पायऱ्या
$\chi \lif \psi$ आणि $\chi$ असतात; म्हणजे $\Gamma \cup \{\varphi\} \Proves \chi \lif \psi$
आणि $\Gamma \cup \{\varphi\} \Proves \chi$. त्यांच्याशी संबंधित
निष्पत्त्या लहान असल्यामुळे त्यांना विगमन गृहीतक लागू होते. म्हणून
आपल्याकडे पुढील दोन्ही आहेत:
\begin{align*}
  & \Gamma \Proves \varphi \lif (\chi \lif \psi); \\
  & \Gamma \Proves \varphi \lif \chi.
\end{align*}
पण \ref{fol:axd:prp:ax:lif2} वरून पुढीलही मिळते:
\[\Gamma \Proves (\varphi \lif (\chi \lif \psi)) \lif
((\varphi\lif \chi)  \lif (\varphi \lif \psi)),
\]
आणि \ref{fol:axd:ded:prop:mp} चे दोन उपयोग केल्यावर आवश्यकतेनुसार
$\Gamma \Proves \varphi \lif \psi$ मिळते.
\end{proof}

निगमन प्रमेय खरे ठरावे यासाठीच \ref{fol:axd:prp:ax:lif1} आणि
\ref{fol:axd:prp:ax:lif2} यांची नेमकी अशी निवड केली होती, हे लक्षात घ्या.

निष्पन्नता विषयीची पुढील काही उपयुक्त तथ्ये सरावासाठी सोडली आहेत.

\begin{prop}
  \label{fol:axd:ded:prop:derivfacts}
\begin{enumerate}
\item $\Proves (\varphi \lif \psi) \lif ((\psi \lif \chi)
  \lif (\varphi \lif \chi))$; \label{fol:axd:ded:derivfacts:a}

\item $\Gamma \cup \{ \lnot \varphi\}
  \Proves \lnot \psi$ असेल, तर $\Gamma \cup \{ \psi\} \Proves
  \varphi$ (प्रतिपरिवर्तन); \label{fol:axd:ded:derivfacts:b}
\item  $\{ \varphi, \lnot\varphi\} \Proves
    \psi$ (Ex Falso Quodlibet, म्हणजे असत्यापासून काहीही; स्फोट); \label{fol:axd:ded:derivfacts:c}
\item  $\{ \lnot\lnot\varphi\} \Proves
  \varphi$ (द्विनकरण विलोपन);\label{fol:axd:ded:derivfacts:d}
\item $\Gamma \Proves \lnot\lnot\varphi$ असेल, तर $\Gamma \Proves
  \varphi$;\label{fol:axd:ded:derivfacts:e}
\end{enumerate}
\end{prop}


\begin{prob}
  \ref{fol:axd:ded:prop:derivfacts} सिद्ध करा.
\end{prob}















\section{संख्यापकांसह निगमन प्रमेय}\label{fol:axd:ddq:sec}

\begin{thm}[निगमन प्रमेय]
\label{fol:axd:ddq:thm:deduction-thm-q} जर $\Gamma \cup \{\varphi\} \Proves \psi$ असेल,
तर $\Gamma \Proves \varphi \lif \psi$.
\end{thm}

\begin{proof}
$\Gamma \cup \{\varphi\}$ पासून $\psi$ च्या निष्पत्तीच्या लांबीवर
आपण पुन्हा विगमन करू.

विगमनाच्या आधारपायरीचा पुरावा~\ref{fol:axd:ded:thm:deduction-thm} च्या
पुराव्यातील आधारपायरीच्या पुराव्यासारखाच आहे.

विगमन पायरीसाठी, $\Gamma \cup \{\varphi\}$ पासून $\psi$ ची निष्पत्ती
एखाद्या निगमन नियमाने समर्थित असलेल्या~$\psi$ या पायरीने संपते असे पुन्हा
समजा. निगमन नियम विध्यात्मक अनुमान असेल, तर आपण
\ref{fol:axd:ded:thm:deduction-thm} च्या पुराव्याप्रमाणेच पुढे जाऊ. निगमन
नियम \QR असेल, तर $\psi \ident \chi \lif \lforall[x][\theta(x)]$ आहे आणि
$\chi \lif \theta(a)$ या रूपाचे सूत्र त्या निष्पत्तीमध्ये आधी
येते, जेथे $a$ हे~$\chi$, $\varphi$, किंवा $\Gamma$ मध्ये येत नाही. म्हणून
आपल्याकडे पुढील आहे:
\begin{align*}
  \Gamma \cup \{\varphi\} & \Proves \chi \lif \theta(a),\\
  \intertext{आणि विगमन गृहीतक लागू होते; म्हणजे पुढील मिळते:}
    \Gamma & \Proves \varphi \lif (\chi \lif \theta(a)).\\
  \intertext{पुढील सूत्रावरून}
  & \Proves (\varphi \lif (\chi \lif \theta(a))) \lif ((\varphi \land \chi) \lif \theta(a))\\
  \intertext{आणि विध्यात्मक अनुमानामुळे पुढील मिळते:}
  \Gamma & \Proves (\varphi \land \chi) \lif \theta(a).\\
  \intertext{आयगेन चराची अट अजूनही लागू असल्यामुळे, या निष्पत्तीमध्ये \QR ने समर्थित एक पायरी जोडून पुढील मिळवता येते:}
    \Gamma & \Proves (\varphi \land \chi) \lif \lforall[x][\theta(x)].\\
    \intertext{आपल्याकडे पुढीलही आहे:}
    & \Proves ((\varphi \land \chi) \lif \lforall[x][\theta(x)]) \lif (\varphi \lif (\chi \lif \lforall[x][\theta(x)])),\\
    \intertext{म्हणून विध्यात्मक अनुमानाने,}
    \Gamma & \Proves \varphi \lif (\chi \lif \lforall[x][\theta(x)]),
\end{align*}

म्हणजे $\Gamma \Proves \varphi \lif \psi$.


$\psi$ ला \QR या नियमाने समर्थन मिळते, पण त्याचे रूप
$\lexists[x][\theta(x)] \lif \chi$ आहे, हे प्रकरण आपण सरावासाठी सोडतो.
\end{proof}

\begin{prob}
  \ref{fol:axd:ddq:thm:deduction-thm-q} चा पुरावा पूर्ण करा.
\end{prob}












\section{निष्पन्नता आणि सुसंगतता}\label{fol:axd:prv:sec}

आता आपण निष्पन्नता संबंधाचे अनेक गुणधर्म सिद्ध करू. ते स्वतंत्रपणेही
महत्त्वाचे आहेत, पण संपूर्णता प्रमेयाच्या सिद्धतेत प्रत्येकाची भूमिका असेल.

\begin{prop}\label{fol:axd:prv:prop:provability-contr}
  जर $\Gamma \Proves \varphi$ आणि $\Gamma \cup \{\varphi\}$ विसंगत असेल, तर
  $\Gamma$ विसंगत आहे.
\end{prop}

\begin{proof}
  जर $\Gamma \cup \{\varphi\}$ विसंगत असेल, तर $\Gamma \cup \{\varphi\}
  \Proves \lfalse$. \ref{fol:axd:ptn:prop:reflexivity} वरून प्रत्येक
  $\psi \in \Gamma$ साठी $\Gamma \Proves \psi$. गृहीतकानुसार
  $\Gamma \Proves \varphi$ सुद्धा असल्यामुळे प्रत्येक $\psi \in \Gamma \cup
  \{\varphi\}$ साठी $\Gamma \Proves \psi$. \ref{fol:axd:ptn:prop:transitivity}
  वरून $\Gamma \Proves \lfalse$, म्हणजे $\Gamma$ विसंगत आहे.
\end{proof}

\begin{prop}
\label{fol:axd:prv:prop:prov-incons}
$\Gamma \Proves \varphi$ तेव्हा आणि केवळ तेव्हाच, जेव्हा $\Gamma \cup \{\lnot \varphi\}$
विसंगत असते.
\end{prop}

\begin{proof}
प्रथम $\Gamma \Proves \varphi$ असे समजा. मग \ref{fol:axd:ptn:prop:monotonicity}
वरून $\Gamma \cup \{\lnot \varphi\} \Proves \varphi$. \ref{fol:axd:ptn:prop:reflexivity}
वरून $\Gamma \cup \{\lnot \varphi\} \Proves \lnot \varphi$. तसेच
\ref{fol:axd:prp:ax:lnot2} वरून $\Proves \lnot \varphi \lif (\varphi \lif \lfalse)$.
म्हणून \ref{fol:axd:ded:prop:mp} चे दोन उपयोग करून $\Gamma \cup \{\lnot
\varphi\} \Proves \lfalse$ मिळते.

आता $\Gamma \cup \{\lnot \varphi\}$ विसंगत आहे असे गृहीत धरा, म्हणजे
$\Gamma \cup \{\lnot \varphi\} \Proves \lfalse$. निगमन प्रमेयावरून $\Gamma
\Proves \lnot \varphi \lif \lfalse$. \ref{fol:axd:prp:ax:lfalse2} वरून $\Gamma
\Proves (\lnot \varphi \lif \lfalse) \lif \lnot\lnot \varphi$, म्हणून
\ref{fol:axd:ded:prop:mp} वरून $\Gamma \Proves \lnot\lnot \varphi$. आपल्याकडे
$\Gamma \Proves \lnot\lnot \varphi \lif \varphi$ हेही असल्यामुळे
(\ref{fol:axd:prp:ax:dne}), \ref{fol:axd:ded:prop:mp} च्या आणखी एका उपयोगाने
$\Gamma \Proves \varphi$ मिळते.
\end{proof}

\begin{prob}
$\Gamma \Proves \lnot \varphi$ तेव्हा आणि केवळ तेव्हाच, जेव्हा $\Gamma \cup \{\varphi\}$
विसंगत असते, हे सिद्ध करा.
\end{prob}

\begin{prop}\label{fol:axd:prv:prop:explicit-inc}
  जर $\Gamma \Proves \varphi$ आणि $\lnot \varphi \in \Gamma$ असेल, तर $\Gamma$
  विसंगत आहे.
\end{prop}

\begin{proof}
  \ref{fol:axd:prp:ax:lnot2} वरून $\Gamma \Proves \lnot \varphi \lif (\varphi \lif \lfalse)$.
  \ref{fol:axd:ded:prop:mp} चे दोन उपयोग करून $\Gamma \Proves \lfalse$ मिळते.
\end{proof}


\begin{prop}\label{fol:axd:prv:prop:provability-exhaustive}
  $\Gamma \cup \{\varphi\}$ आणि $\Gamma \cup \{\lnot \varphi\}$ हे दोन्ही संच
  विसंगत असतील, तर $\Gamma$ विसंगत आहे.
\end{prop}

\begin{proof}
सराव.
\end{proof}


\begin{prob}
  \ref{fol:axd:prv:prop:provability-exhaustive} सिद्ध करा
\end{prob}
















\section{निष्पन्नता आणि विधानीय संयोजक}\label{fol:axd:ppr:sec}

\begin{explain}
  स्वयंसिद्धकीय निगमनाचा निष्पन्नता संबंध~$\Proves$ हा विधानीय
  संयोजकांशी संबंधित काही मूलभूत तथ्ये सिद्ध करण्यासाठी पुरेसा सामर्थ्यवान
  आहे; उदाहरणार्थ, $\varphi \land \psi \Proves \varphi$ आणि $\varphi, \varphi \lif \psi
  \Proves \psi$ (विध्यात्मक अनुमान). ही तथ्ये संपूर्णता प्रमेयाच्या
  सिद्धतेसाठी आवश्यक आहेत.
\end{explain}

\begin{prop}\label{fol:axd:ppr:prop:provability-land}
  \begin{enumerate}
  \item \label{fol:axd:ppr:prop:provability-land-left} $\varphi \land \psi \Proves
    \varphi$ आणि $\varphi \land \psi \Proves \psi$ ही दोन्ही तथ्ये सत्य आहेत.
  \item \label{fol:axd:ppr:prop:provability-land-right} $\varphi, \psi \Proves \varphi \land \psi$.
  \end{enumerate}
\end{prop}

\begin{proof}
  \begin{enumerate}
    \item \ref{fol:axd:prp:ax:land1} आणि \ref{fol:axd:prp:ax:land2} यांपासून,
      अनुक्रमे विध्यात्मक अनुमान लावून.

  \item \ref{fol:axd:prp:ax:land3} पासून विध्यात्मक अनुमान दोनदा लावून.
  \end{enumerate}
\end{proof}


\begin{prop}\label{fol:axd:ppr:prop:provability-lor}
  \begin{enumerate}
  \item $\varphi \lor \psi, \lnot \varphi, \lnot \psi$ विसंगत आहे.
  \item $\varphi \Proves \varphi \lor \psi$ आणि $\psi \Proves \varphi \lor \psi$ ही दोन्ही
    तथ्ये सत्य आहेत.
  \end{enumerate}
\end{prop}

\begin{proof}
  \begin{enumerate}
  \item \ref{fol:axd:prp:ax:lnot2} पासून $\Proves \lnot \varphi \lif (\varphi
    \lif \lfalse)$ आणि $\Proves \lnot \psi \lif (\psi \lif \lfalse)$
    मिळतात. म्हणून निगमन प्रमेयाने $\{\lnot \varphi\} \Proves \varphi \lif
    \lfalse$ आणि $\{\lnot \psi\} \Proves \psi \lif \lfalse$ मिळतात.
    \ref{fol:axd:prp:ax:lor3} पासून $\{\lnot \varphi, \lnot \psi\} \Proves (\varphi
    \lor \psi) \lif \lfalse$ मिळते. निगमन प्रमेयाने $\{\varphi \lor \psi,
    \lnot \varphi, \lnot \psi\} \Proves \lfalse$.

  \item \ref{fol:axd:prp:ax:lor1} आणि \ref{fol:axd:prp:ax:lor2} यांपासून विध्यात्मक
    अनुमानाने.
  \end{enumerate}
\end{proof}

\begin{prop}\label{fol:axd:ppr:prop:provability-lif}
  \begin{enumerate}
  \item \label{fol:axd:ppr:prop:provability-lif-left} $\varphi, \varphi \lif \psi \Proves \psi$.
  \item \label{fol:axd:ppr:prop:provability-lif-right}
    $\lnot \varphi \Proves \varphi \lif \psi$ आणि $\psi \Proves \varphi \lif \psi$ ही दोन्ही
    तथ्ये सत्य आहेत.
  \end{enumerate}
\end{prop}

\begin{proof}
  \begin{enumerate}
  \item आपण पुढील निष्पन्न करू शकतो:
    \begin{derivation}
      1. & $\varphi$ & \Hyp\\
      2. & $\varphi \lif \psi$ & \Hyp\\
      3. & $\psi$ & 1, 2, \MP
    \end{derivation}

  \item अनुक्रमे \ref{fol:axd:prp:ax:lnot2}, \ref{fol:axd:prp:ax:lif1} आणि निगमन
    प्रमेय वापरून.
  \end{enumerate}
\end{proof}















\section{निष्पन्नता आणि संख्यापक}\label{fol:axd:qpr:sec}

\begin{explain}
  संपूर्णता प्रमेयासाठी या विभागात स्थापित केलेली~$\Proves$ संबंधी तथ्ये
  स्वयंसिद्धकीय निगमनातून निष्पन्न होतात, हेही आवश्यक आहे.
\end{explain}

\begin{thm}
\label{fol:axd:qpr:thm:strong-generalization} जर $c$ हे $\Gamma$ किंवा $\varphi(x)$ मध्ये
न येणारे स्थिरांक असेल आणि $\Gamma \Proves \varphi(c)$, तर $\Gamma
\Proves \lforall[x][\varphi(x)]$.
\end{thm}

\begin{proof}
निगमन प्रमेयाने $\Gamma \Proves \ltrue \lif \varphi(c)$. $c$ हे
$\Gamma$ किंवा~$\top$ मध्ये येत नसल्यामुळे $\Gamma \Proves
\ltrue \lif \lforall[x][\varphi(x)]$ मिळते. त्यानंतर सत्य स्थिरांकाचे
स्वयंसिद्धक आणि विध्यात्मक अनुमान वापरून $\Gamma \Proves
\lforall[x][\varphi(x)]$ मिळते.

\end{proof}

\begin{prop}
\label{fol:axd:qpr:prop:provability-quantifiers}
\begin{enumerate}
\item $\varphi(t) \Proves \lexists[x][\varphi(x)]$.

\item $\lforall[x][\varphi(x)] \Proves \varphi(t)$.
\end{enumerate}
\end{prop}

\begin{proof}
\begin{enumerate}
\item \ref{fol:axd:qua:ax:q2} आणि निगमन प्रमेय वापरून.
\item \ref{fol:axd:qua:ax:q1} आणि निगमन प्रमेय वापरून.
\end{enumerate}
\end{proof}












\section{निर्दोषता}\label{fol:axd:sou:sec}

\begin{explain}
स्वयंसिद्धकीय निगमनासारखी निष्पत्ती-पद्धत
\emph{निर्दोष} असते, जर ती प्रत्यक्षात निष्पन्न न होणाऱ्या गोष्टी
निष्पन्न करू शकत नसेल. त्यामुळे निर्दोषता हा निष्पत्ती-पद्धतींसाठी
हमी दिलेल्या सुरक्षिततेचा एक प्रकार आहे. कोणता सिद्धता-उपपत्तीय गुणधर्म
विचारात आहे यावर अवलंबून, उदाहरणार्थ, आपल्याला पुढील गोष्टी हव्या असतात:
\begin{enumerate}
\item प्रत्येक निष्पन्न करता येण्याजोगे~$\varphi$ वैध असते;
\item $\varphi$ हे इतर काही सूत्रे~$\Gamma$ यांच्यापासून निष्पन्न करता येण्याजोगे असेल,
  तर ते त्यांचा तार्किक परिणामही असते;
\item सूत्रांचा संच~$\Gamma$ विसंगत असेल, तर तो अपूर्ततायोग्य असतो.
\end{enumerate}
हे निष्पत्ती-पद्धतीचे महत्त्वाचे गुणधर्म आहेत. यांपैकी एखादा गुणधर्म
खरा नसेल, तर निष्पत्ती-पद्धतीत दोष आहे---ती खूप जास्त गोष्टी
निष्पन्न करेल. म्हणून निष्पत्ती-पद्धतीची निर्दोषता सिद्ध करणे अत्यंत
महत्त्वाचे आहे.
\end{explain}

\begin{prop}
  जर $\varphi$ हे स्वयंसिद्धक असेल, तर
  प्रत्येक संरचना~$\Struct{M}$ आणि चल-विन्यास~$s$ साठी $\Sat{M}{\varphi}[s]$.
\end{prop}

\begin{proof}
सर्व स्वयंसिद्धके वैध आहेत हे तपासावे लागते. उदाहरणार्थ,
  \ref{fol:axd:qua:ax:q1} चे प्रकरण पाहू: $t$ हे $\varphi$ मध्ये $x$ साठी
  आदेशनासाठी मुक्त आहे असे समजा आणि
  $\Sat{M}{\lforall[x][\varphi]}[s]$ असे गृहीत धरा. मग पूर्ततेच्या व्याख्येनुसार
  प्रत्येक $\varAssign{s'}{s}{x}$ साठी $\Sat{M}{\varphi}[s']$, आणि विशेषतः
  $s'(x) = \Value{t}{M}[s]$ असताना हे खरे आहे.
  \ref{fol:syn:ext:prop:ext-formulas} नुसार
  $\Sat{M}{\Subst{\varphi}{t}{x}}[s]$. यावरून
  $\Sat{M}{(\lforall[x][\varphi] \lif \Subst{\varphi}{t}{x})}[s]$ हे सिद्ध होते.
\end{proof}

\begin{thm}[निर्दोषता]
\label{fol:axd:sou:thm:soundness}
जर $\Gamma \Proves \varphi$, तर $\Gamma \Entails \varphi$.
\end{thm}

\begin{proof}
$\Gamma$ पासून $\varphi$ च्या निष्पत्तीच्या लांबीवर विगमन करू. अनुमानांनी
समर्थित कोणत्याही पायऱ्या नसतील, तर निष्पत्तीमधील सर्व सूत्रे
ही स्वयंसिद्धकांची दृष्टांते असतात किंवा~$\Gamma$ मध्ये असतात. आधीच्या
विधानानुसार सर्व स्वयंसिद्धके वैध आहेत;
म्हणून $\varphi$ हे स्वयंसिद्धक असेल, तर $\Gamma \Entails \varphi$. जर
$\varphi \in \Gamma$, तर स्पष्टपणे $\Gamma \Entails \varphi$.

जर~$\varphi$ च्या निष्पत्तीमधील शेवटची पायरी विध्यात्मक अनुमानाने समर्थित
असेल, तर त्या निष्पत्तीमध्ये $\psi$ आणि $\psi \lif \varphi$ ही सूत्रे
असतात; आणि त्या सूत्रे वर संपणाऱ्या निष्पत्तीच्या भागांना विगमन
गृहीतक लागू होते (कारण त्यांच्यात अनुमानाने समर्थित किमान एक पायरी कमी
असते). म्हणून विगमन गृहीतकानुसार $\Gamma \Entails \psi$ आणि
$\Gamma \Entails \psi \lif \varphi$. मग
\ref{fol:syn:sem:thm:sem-deduction}
नुसार $\Gamma \Entails \varphi$.

आता शेवटची पायरी \QR ने समर्थित आहे असे समजा. मग त्या पायरीचे
रूप $\chi \lif \lforall[x][\psi(x)]$ असते आणि तिच्याआधीची पायरी
$\chi \lif \psi(c)$ असते; येथे $c$ हे $\Gamma$, $\chi$ किंवा
$\lforall[x][B(x)]$ यांमध्ये येत नाही. विगमन गृहीतकानुसार
$\Gamma \Entails \chi \lif \psi(c)$.
\ref{fol:syn:sem:thm:sem-deduction} नुसार $\Gamma \cup \{\chi\} \Entails
\psi(c)$.


$\Sat{M}{\Gamma \cup
  \{\chi\}}$ अशी एखादी रचना~$\Struct{M}$ विचारात घ्या. आपल्याला
$\Sat{M}{\lforall[x][\psi(x)]}$ हे दाखवायचे आहे.
$\lforall[x][\psi(x)]$ हे वाक्य असल्यामुळे याचा अर्थ प्रत्येक
चल-विन्यास~$s$ साठी $\Sat{M}{\psi(x)}[s]$ दाखवायचे आहे
(\ref{fol:syn:ass:prop:sat-quant}). $\Gamma \cup \{\chi\}$ मध्ये केवळ
वाक्ये असल्यामुळे \ref{fol:syn:sat:defn:satisfaction} नुसार प्रत्येक
$\theta \in \Gamma$ साठी $\Sat{M}{\theta}[s]$. $\Struct{M'}$ ही
$\Struct{M}$ सारखीच रचना असू द्या, फक्त $\Assign{c}{M'} = s(x)$.
$c$ हे~$\Gamma$ किंवा~$\chi$ मध्ये येत नसल्यामुळे
\ref{fol:syn:ext:cor:extensionality-sent} नुसार
$\Sat{M'}{\Gamma \cup \{\chi\}}$. $\Gamma \cup \{\chi\}
\Entails \psi(c)$ असल्यामुळे $\Sat{M'}{\psi(c)}$. $\psi(c)$ हे वाक्य
असल्यामुळे \ref{fol:syn:ass:prop:sentence-sat-true} नुसार
$\Sat{M'}{\psi(c)}[s]$. \ref{fol:syn:ext:prop:ext-formulas} नुसार
$\Sat{M'}{\psi(x)}[s]$ तेव्हा आणि केवळ तेव्हाच, जेव्हा
$\Sat{M'}{\psi(c)}[s]$ (लक्षात घ्या की $\psi(c)$ म्हणजे
$\Subst{\psi(x)}{c}{x}$). म्हणून $\Sat{M'}{\psi(x)}[s]$. $c$ हे~$\psi(x)$
मध्ये येत नसल्यामुळे \ref{fol:syn:ext:prop:extensionality} नुसार
$\Sat{M}{\psi(x)}[s]$. पण $s$ हा स्वेच्छ चल-विन्यास होता, म्हणून
$\Sat{M}{\lforall[x][\psi(x)]}$. त्यामुळे $\Gamma \cup \{\chi\} \Entails
\lforall[x][\psi(x)]$. \ref{fol:syn:sem:thm:sem-deduction} नुसार
$\Gamma \Entails \chi \lif \lforall[x][\psi(x)]$.


$\varphi$ हे \QR ने समर्थित पण $\lexists[x][\psi(x)] \lif \chi$ या रूपाचे असलेले
प्रकरण सरावासाठी सोडले आहे.
\end{proof}


\begin{prob}
  \ref{fol:axd:sou:thm:soundness} ची सिद्धता पूर्ण करा.
\end{prob}


\begin{cor}
\label{fol:axd:sou:cor:weak-soundness}
जर $\Proves \varphi$, तर $\varphi$ हे वैध आहे.
\end{cor}

\begin{cor}
\label{fol:axd:sou:cor:consistency-soundness}
जर $\Gamma$ पूर्ततायोग्य असेल, तर तो सुसंगत आहे.
\end{cor}

\begin{proof}
प्रतिलोम विधान सिद्ध करू. $\Gamma$ सुसंगत नाही असे समजा. मग
$\Gamma \Proves \lfalse$, म्हणजे~$\Gamma$ पासून $\lfalse$ ची
निष्पत्ती असते. \ref{fol:axd:sou:thm:soundness} नुसार $\Gamma$ पूर्ण करणाऱ्या
प्रत्येक
संरचना~$\Struct{M}$
ने~$\lfalse$ पूर्ण केले पाहिजे. प्रत्येक
संरचना~$\Struct{M}$ साठी
  $\Sat/{M}{\lfalse}$ असल्यामुळे कोणतीही
$\Struct{M}$ $\Gamma$ पूर्ण करू शकत नाही;
म्हणजे $\Gamma$ अपूर्ततायोग्य आहे.
\end{proof}














\section{निष्पत्ती सह एकरूपता}\label{fol:axd:ide:sec}

निष्पत्त्यांमध्ये $\eq$ समाविष्ट करण्यासाठी आपण फक्त नवीन
स्वयंसिद्धक-आकार जोडतो. निष्पत्ती आणि $\Proves$ यांच्या व्याख्या
तशाच राहतात; आता आपण नवीन स्वयंसिद्धकेही मान्य करतो.

\begin{defn}[एकरूपता साठीची स्वयंसिद्धके]
\begin{align}
\label{fol:axd:ide:ax:id1} & \eq[t][t], \\
\label{fol:axd:ide:ax:id2} & \eq[t_1][t_2] \lif (\psi(t_1) \lif
  \psi(t_2)),
\end{align}
येथे $t$, $t_1$, $t_2$ ही कोणतीही बंद पदे आहेत.
\end{defn}

\begin{prop}\label{fol:axd:ide:prop:sound}
  \ref{fol:axd:ide:ax:id1} आणि \ref{fol:axd:ide:ax:id2} ही स्वयंसिद्धके वैध आहेत.
\end{prop}

\begin{proof}
  सराव.
\end{proof}

\begin{prob}
\ref{fol:axd:ide:prop:sound} सिद्ध करा.
\end{prob}

\begin{prop}\label{fol:axd:ide:prop:iden1}
 कोणतेही पद $t$ आणि कोणताही संच~$\Gamma$ यांसाठी
 $\Gamma \Proves \eq[t][t]$.
\end{prop}

\begin{prop}\label{fol:axd:ide:prop:iden2}
  जर $\Gamma \Proves \varphi(t_1)$ आणि $\Gamma \Proves
  \eq[t_1][t_2]$, तर $\Gamma \Proves \varphi(t_2)$.
\end{prop}

\begin{proof}
पुढील सूत्र
\[(\eq[t_1][t_2] \lif (\varphi(t_1) \lif \varphi(t_2)))
\]
हा~\ref{fol:axd:ide:ax:id2} चा दृष्टांत आहे. निष्कर्ष~\MP चे दोन उपयोग केल्यावर
मिळतो.
\end{proof}



